We investigated whether AMR-structured rewriting can reduce the score of a semantic watermark detector without collapsing output quality. The results separate detector vulnerability from usable attack success. Boundary-only rewriting produces the strongest raw escape rate (93.3\%), but that number is not a suitable headline because it comes with a clear quality cost. Under the stricter conjunction of detector escape and validity preservation, the robustness-backed configuration succeeds on 9 of 30 paragraphs, yielding 30.0\% strict-valid ASR with an exact 95\% confidence interval of [14.7\%, 49.4\%].

Two checks support that result. The strict value stays unchanged across the 0.65--0.80 search-threshold plateau, and all 120 stored candidates retain their binary detector decisions over 50 recomputations. The candidate-budget frontier also shows that more search is not automatically better: the best observed strict result is 33.3\% at $K=1$, while larger budgets reduce detector scores without monotonically improving semantic validity. The task is therefore a constrained, multi-objective search, not pure score minimization.

These experiments show that, under one RealNews-based subset and one fixed AMR detector configuration, structure-aware rewriting can produce a limited set of paragraphs that pass both the detector and the project's validity filter. They do not show a general break of semantic watermarking, nor do they establish transfer to other detectors, corpora, or generation settings. The defensible takeaway is that AMR structure is a meaningful evaluation surface and a concrete attack surface in this pipeline.
